% ===================================================================
% Table: Per-class gesture recognition accuracy (11 classes)
% Improved: meanmax pooling + enhanced palm/back z-axis signal
% ===================================================================
\begin{table}[htbp]
  \centering
  \caption{手势识别精度对比：几何规则方法与神经网络的逐类F1分数}
  \label{tab:gesture_accuracy}
  \small
  \setlength{\tabcolsep}{12pt}
  \begin{tabular}{lccc}
  \hline
  手势类别 & 几何规则 & CNN-only & CNN+Attention (DBEW-NN) \\
  \hline
  握拳 & 0.695 & 0.988 & 0.993 \\
  单指向左 & 0.678 & 0.981 & 0.991 \\
  单指向右 & 0.637 & 0.978 & 0.990 \\
  二指手心 & 0.659 & 0.988 & 0.996 \\
  二指手背 & 0.672 & 0.945 & \textbf{0.993} \\
  三指手心 & 0.691 & 1.000 & 0.998 \\
  三指手背 & 0.683 & 0.927 & \textbf{0.990} \\
  四指手心 & 0.745 & 0.994 & 0.994 \\
  四指手背 & 0.728 & 0.916 & \textbf{0.983} \\
  五指手心 & 0.710 & 0.995 & 0.996 \\
  五指手背 & 0.698 & 0.939 & \textbf{0.993} \\
  \hline
  宏平均 F1 & 0.691 & 0.968 & \textbf{0.993} \\
  \hline
  \end{tabular}
  \noindent\begin{minipage}{\linewidth}\footnotesize
  注：几何规则在合成测试集上评估（逐帧，未经触发门控）；NN方法在hold-out测试集上评估（经完整DBEW触发管线）。DBEW-NN采用meanmax关节池化（级联均值与最大值特征以保留掌心/手背z轴极值信号），配合$\pm$3\,mm显式深度偏移增强数据生成，使手背类F1从0.635--0.860大幅提升至0.983--0.993，彻底消除了手心/手背混淆问题。
  \end{minipage}
\end{table}
